% Ad Familiares 12.30 (ad-familiares-12-30) --- English translation
% Translated by Claude (Anthropic) from the Latin (Perseus).
% Style guide: STYLE.md.

\ciceroLetterOpener{CICERO CORNIFICIO S.}

\ciceroSection{1} What, really? No one comes to you with a letter from me except litigants? Plenty of those, to be sure --- for you yourself have made it the case that no one thinks himself recommended to you without a letter of mine. But which of your people ever told me there was someone to whom I could entrust one, that I did not entrust it? And what is more delightful to me than --- since I cannot speak with you face to face --- either to write to you or to read your letters? What more often vexes me is the load of business that hampers me, so that no chance is given me of writing to you at my own discretion. I would be pelting you not with letters but with whole volumes; by which indeed I ought to be provoked by you. For however busy you are, you have more leisure than I do; or, if not even you have any leisure, do not be shameless or vex me by demanding more frequent letters from me when you yourself send rarely.

\ciceroSection{2} For where before I was pulled apart by the greatest occupations --- because I reckoned the Republic must be guarded by me with every care --- now at this time I am pulled apart far more violently. For just as those are taken more gravely ill who, when they seem to have got past the disease, fall back into it afresh, so we are labouring more violently, who, with the war beaten down and almost done away with, are now trying to wage it renewed. But enough of this.

\ciceroSection{3} Persuade yourself, my dear Cornificius, that I am not of so weak a spirit --- not to say so inhuman --- that you can outdo me either in services or in affection. I did not doubt it, in fact, but Chaerippus has made your affection toward me much better known to me. What a man! Always congenial to me, but now actually charming! By Hercules, he has set out your whole face before me, has carried over to me not only your mind and your words. So do not be afraid that I am offended because you wrote to me by the same draft you sent to others. I did indeed ask for a letter to me alone, but not vehemently --- with affection.

\ciceroSection{4} On the expense which you say you are incurring and have incurred for military matters, I can really be no help to you, because the Senate has been left orphaned by the loss of the consuls, and the straits of the public purse are unbelievable: money is being scraped together from every quarter, that what was promised to the soldiers who have served best may be paid out --- and I do not think this can be done without a special levy.

\ciceroSection{5} About Attius Dionysius, I take it there is nothing in it, since Tratorius has said nothing to me. About P. Lucceius I yield to you in nothing --- to make you any more zealous than I am; for the man is intimate with us. But when I pressed the arbiters about a postponement of the date, they showed me they were prevented from granting it both by the formal compromise and by their oath. So Lucceius must come, I think; though, if he has obeyed my letter, when you read this he ought to be at Rome.

\ciceroSection{6} On the other matters, and most of all on the money, you wrote --- while you did not know that Pansa was dead --- what you reckoned you could obtain from him through us. These things would not have failed you, if he were alive; for he held you dear. But after his death we did not see what could be done. About Venuleius, Latinus, and Horatius I warmly approve; the one thing I do not altogether endorse is what you write --- that, to make them bear it with a more level mind, you stripped your own legates also of their lictors. (Men worthy of honour were not to be compared with men worthy of disgrace.) And these, if they do not depart by force of the senatorial decree, I think must be compelled to depart. So much, more or less, in answer to those letters, two copies of which, identical, he brought me at the same time; for the rest, I would have you persuade yourself that my own standing is no dearer to me than yours.
